\documentclass[11pt,a4paper]{article}
\usepackage[T1]{fontenc}
\usepackage[utf8]{inputenc}
\usepackage{lmodern}
\usepackage[margin=26mm,headheight=15pt]{geometry}
\usepackage{amsmath,amssymb,amsthm,mathtools}
\usepackage{microtype,booktabs,longtable,array,enumitem}
\usepackage[dvipsnames]{xcolor}
\usepackage{fancyhdr,xurl}
\usepackage[colorlinks=true,linkcolor=MidnightBlue,citecolor=MidnightBlue,urlcolor=MidnightBlue]{hyperref}
\hypersetup{pdftitle={Definable Surreal Numbers and Omnific Integers},pdfauthor={Research article prepared for the Surreal project},pdfsubject={Ambient definability, HOD, reversible omnific coding, and a transcendence dichotomy}}
\pagestyle{fancy}\fancyhf{}
\fancyhead[L]{\small Definable surreals and omnific integers}
\fancyhead[R]{\small\thepage}
\renewcommand{\headrulewidth}{0.3pt}
\setlength{\parskip}{3pt}
\setlength{\emergencystretch}{2em}
\setlist{itemsep=3pt,topsep=4pt,leftmargin=2em}
\numberwithin{equation}{section}
\allowdisplaybreaks[2]
\newtheorem{theorem}{Theorem}[section]
\newtheorem{lemma}[theorem]{Lemma}
\newtheorem{proposition}[theorem]{Proposition}
\newtheorem{corollary}[theorem]{Corollary}
\theoremstyle{definition}
\newtheorem{definition}[theorem]{Definition}
\newtheorem{example}[theorem]{Example}
\newtheorem{question}[theorem]{Question}
\theoremstyle{remark}
\newtheorem{remark}[theorem]{Remark}
\newcommand{\No}{\mathbf{No}}
\newcommand{\Oz}{\mathbf{Oz}}
\newcommand{\Ord}{\mathrm{Ord}}
\newcommand{\OD}{\mathrm{OD}}
\newcommand{\HOD}{\mathrm{HOD}}
\newcommand{\ZFC}{\mathrm{ZFC}}
\newcommand{\N}{\mathbb N}
\newcommand{\Z}{\mathbb Z}
\newcommand{\Q}{\mathbb Q}
\newcommand{\R}{\mathbb R}
\newcommand{\C}{\mathbb C}
\newcommand{\bd}{\operatorname{b}}
\newcommand{\ct}{\operatorname{ct}}
\newcommand{\supp}{\operatorname{supp}}
\newcommand{\NF}{\operatorname{NF}}
\newcommand{\dcl}{\operatorname{dcl}}
\newcommand{\acl}{\operatorname{acl}}
\newcommand{\Frac}{\operatorname{Frac}}
\newcommand{\TC}{\operatorname{TC}}
\newcommand{\dom}{\operatorname{dom}}
\newcommand{\ran}{\operatorname{ran}}
\newcommand{\sgn}{\operatorname{sgn}}
\newcommand{\Sat}{\operatorname{Sat}}
\newcommand{\Code}{\mathsf C}
\newcommand{\Decode}{\mathsf D}
\newcommand{\Cell}{\mathcal U}
\newcommand{\KH}{K_{\HOD}}
\newcommand{\IH}{I_{\HOD}}
\newcommand{\repoSHA}{\texttt{9a385d3957bdfe3d9ea79f9a524751c90bd2c894}}
\newcommand{\repopath}[1]{\nolinkurl{#1}}
\newcolumntype{P}[1]{>{\raggedright\arraybackslash}p{#1}}
\title{\textbf{Definable Surreal Numbers\\and Omnific Integers}\\[5pt]
\large Ambient definability, reversible coding,\newline
hereditary ordinal definability, and a transcendence dichotomy}
\author{Research article prepared for the Surreal project}
\date{September 23, 2026}
\begin{document}
\maketitle
\begin{abstract}
A surreal number cannot be called definable without specifying a language,
a universe, and permitted parameters. We develop ambient set-theoretic
notions, distinguish them from definability in the ordered field, and give
an algebraic theory of the resulting numbers. For a fixed parameter set,
the ambiently definable surreals form a real closed field; the definable
omnific integers form an integer part whose fraction field is exactly that
field. Ordinal-definable surreals admit the particularly concrete description
$\KH=\No\cap\HOD=\No^{\HOD}$, using canonical sign sequences.

Our main constructions go beyond these closure facts. A parameter-free,
reversible normal-form code sends every surreal into the fixed omnific
interval $\omega<z<\omega+\omega^{1/2}$, with leading term $\omega$,
zero constant term, and coefficients in $\{1,2,3\}$. Consequently,
ordinal definability of all integers in this one interval is equivalent to
$V=\HOD$; the analogous parameter-free condition in a model is equivalent
to pointwise definability of the entire model. More strongly, if
$V\ne\HOD$, the same interval contains an explicitly given
$\Ord$-indexed algebraically independent family over $\KH$. Thus the
failure of ordinal definability is never merely a finite or set-bounded
transcendence defect.

We also determine the first birthday of a non-ordinal-definable surreal,
exhibit failures of termwise definability for Hahn sums, construct
complexity-bounded definability escape ordinals, and show how homogeneous
forcing can create the full transcendence phenomenon without adding any
real numbers. Proofs, scope restrictions, source comparisons, and finite
regression checks are included. The coding and transcendence results are
proposed as original contributions; priority is not established, and the
article is not refereed or formally verified.
\end{abstract}
\vspace{3pt}
\noindent\textbf{Keywords.} Surreal numbers; omnific integers; definability;
HOD; normal forms; algebraic independence; pointwise definable models.
\clearpage
\setcounter{tocdepth}{2}
\begingroup\small\setlength{\parskip}{0pt}
\tableofcontents
\endgroup
\clearpage

\section{The question and the contribution}
\label{sec:intro}

The familiar phrase ``definable real number'' conceals several choices.
A positive root of a rational polynomial is definable in an ordered field;
a computable real can instead be defined by a statement about its
algorithm; and a set-theoretic definition can refer to ordinals, models,
and forcing extensions. These are different notions, not alternative
notations for one fixed countable field. The user-specified reference on
definable reals is a useful starting point \cite{WikiDef}; the foundational
analysis and the pointwise-definability results used below come from
Hamkins--Linetsky--Reitz \cite{HLR}.

For surreals the distinction becomes essential. The pure ordered field
has no parameter-free definable infinite element, whereas the first
infinite ordinal $\omega$ has a very short parameter-free definition in
set theory. An omnific integer can simultaneously have elementary-looking
normal-form coefficients and contain the information of an arbitrary
set of ordinals. Our goal is to make these facts into a coherent theory
rather than to announce an unspecified collection of ``describable''
numbers.

\subsection{Three central results}

Let $\No$ be represented by ordinal-length sign sequences, and let $\Oz$
be the classical omnific-integer ring. Write
\begin{equation}\label{eq:maincell}
 \Cell=\{z\in\Oz:\ \omega<z<\omega+\omega^{1/2},\ \ct(z)=0\}.
\end{equation}
The exponentiation $x\mapsto\omega^x$ throughout this article is the
\emph{Conway monomial map}, not $x\mapsto\exp(x\log\omega)$.
The main results are the following.

\paragraph{Reversible compression.}
There is a parameter-free set-theoretically definable injection
$\Code:\No\longrightarrow\Cell$ with a parameter-free definable inverse
on its image. For every permitted parameter family, a number and its
code have exactly the same ambient definability degree
(Theorem~\ref{thm:code}). This is not asserted to be a ring homomorphism
or to be definable in the pure ordered field.

\paragraph{A one-interval test for the universe.}
Every element of $\Cell$ is ordinal definable if and only if $V=\HOD$.
In the external semantics of a model $M\models\ZFC$, every element of
$\Cell^M$ is parameter-free definable in $M$ if and only if $M$ is
pointwise definable (Theorems~\ref{thm:globalOD} and
\ref{thm:pointwise}). Thus a fixed leading asymptotic term does not put
a bound on set-theoretic information.

\paragraph{A transcendence dichotomy.}
Either $\No=\KH$, or $\Cell$ contains an $\Ord$-indexed algebraically
independent family over $\KH$ (Theorem~\ref{thm:dichotomy}). The proof
is explicit: starting with any $t\in(0,1)\setminus\KH$, use
\begin{equation}\label{eq:introfam}
 z_\alpha=\omega+\omega^{\,t\omega^{-(\alpha+1)}}
 \qquad(\alpha\in\Ord).
\end{equation}
The obstruction to a polynomial relation is separation of normal-form
supports into distinct additive cosets of $\KH$.

\subsection{What is prior work, and what is proposed here}

The sign construction, normal forms, the monomial map, real closedness,
and the omnific integer part are classical \cite{Conway,Gonshor,vdDE,LM}.
Quantifier elimination, ordinal definability, HOD, partial truth
predicates, and Mostowski collapse are standard tools
\cite{Marker,Jech,HLR}. We do not claim these tools, or their routine
closure consequences, as new.

Chen, Hamkins, and Yang have announced that the surreal ordered field
expanded by birthday order is bi-interpretable with set theory
\cite{CHY}. The repository explicitly credits that announcement in its
birthday-cutoff report. In particular, the general idea of encoding
well-founded membership graphs by surreal signs is \emph{not} a new
contribution of this article. No theorem below needs an unverified
version of that announced bi-interpretation.

The candidate-original contribution is the particular reversible
\emph{omnific} code with fixed leading data, its one-interval
pointwise-definability criterion, and especially the explicit
\emph{proper-class transcendence dichotomy in that same interval}.
The birthday and complexity results are useful refinements and
applications, with less ambitious originality claims. The literature
search did not locate these precise statements. It was not exhaustive;
no theorem is described as settling a named open problem in the
published literature.

\subsection{Repository snapshot and dependency boundary}

The repository was inspected at commit
\begin{center}\small\repoSHA.\end{center}
The source catalogue and introductions to the birthday-cutoff and
omnific-Diophantine reports were read through the GitHub connector
\cite{Repo,RepoBirthday,RepoOmnific}. The catalogue describes AI-assisted
research drafts and separates written arguments from Lean coverage.
This paper follows that distinction. It imports no nonclassical
Diophantine, quotient, automorphism, or bi-interpretation claim from
those drafts. Appendix~\ref{app:audit} records the limited audit scope.

\section{Foundations and canonical representations}
\label{sec:foundations}

\subsection{Universes, models, and class notation}

Our mathematical statements are made in ZFC, with proper-class notation
understood schematically. Each individual surreal is a \emph{set}.
A class function used below is specified by one set-theoretic formula;
all its values and all sums at which it is evaluated are set-sized.
An $\Ord$-indexed independent family means a definable class assignment
whose every finite subfamily is algebraically independent. It does not
mean that the entire family is a set or that its class sum exists.

When a literal satisfaction relation is needed, fix a transitive
set model $M\models\ZFC$ in an ambient metatheory, and interpret
$\No^M$, $\Oz^M$, and the relevant definitions internally. This
conditional framework does not assert that ZFC proves the existence
of such an $M$. More generally, the model-theoretic statements explicitly
identified as external apply to arbitrary set models. For nonstandard
models, their ``ordinals,'' ``finite sequences,'' and surreal codes
are internal objects, not necessarily actual well-founded sequences.

\subsection{Signs and normal forms}

A canonical surreal is a function
\[
 s:\alpha\longrightarrow\{-,+\},\qquad\alpha\in\Ord.
\]
Its birthday is $\bd(s)=\alpha$. Comparison is by the first disagreement,
with a missing sign ordered between $-$ and $+$. The ordinal $\alpha$
itself is represented by the all-plus sequence of length $\alpha$.
The standard real embedding and the arithmetic operations are
parameter-free definable in set theory \cite{Conway,Gonshor}.

Every nonzero surreal has a unique Conway normal form
\begin{equation}\label{eq:normal}
 x=\sum_{\xi<\lambda}r_\xi\omega^{g_\xi},
 \qquad r_\xi\in\R\setminus\{0\},\quad
 g_\xi>g_\eta\ \ (\xi<\eta<\lambda).
\end{equation}
The length $\lambda$ is an ordinal, and the support
$\supp(x)=\{g_\xi:\xi<\lambda\}$ is reverse well-ordered and set-sized.
Conversely, any such coefficient/exponent sequence has a surreal sum.
Addition and multiplication are Hahn addition and convolution. We use
$\NF(x)$ for the \emph{whole indexed sequence}, including its length,
not merely the collection of its individual entries.

The uniqueness theorem supplies parameter-free set-theoretic operations
$x\mapsto\NF(x)$ and evaluation of valid normal forms. It supplies no
algorithm for arbitrary infinite codes and no assertion of topological
convergence of finite partial sums. The primary factorisation paper
\cite[Sections 1--2]{LM} is a convenient modern reference for the
normal-form convention and the distinction between series fields and
omnific subrings.

\begin{lemma}[Ordinal domination]\label{lem:ordinalbound}
Every set $S\subseteq\No$ is strictly bounded in absolute value by a
surreal ordinal. Moreover, the least ordinal $\delta$ satisfying
$|s|<\delta$ for every $s\in S$ is definable from $S$.
\end{lemma}
\begin{proof}
A sign sequence of length $\alpha$ lies between the all-minus and
all-plus sequences of length $\alpha$, hence $|s|\leq\bd(s)$.
Replacement gives the set of birthdays. An ordinal strictly exceeding
all of them gives the required bound. The qualifying ordinals form a
nonempty definable class, so it has a least member, uniquely specified
from $S$.
\end{proof}

\subsection{Omnific integers and the correct floor formula}

\begin{definition}
An \emph{omnific integer} is a surreal whose normal form has no negative
exponents and whose coefficient at exponent $0$ is an ordinary integer.
Thus
\begin{equation}\label{eq:Oz}
 \Oz=\Pi\oplus\Z,\qquad
 \Pi=\{x\in\No:\supp(x)\subseteq\No_{>0}\}.
\end{equation}
The empty support is permitted, so $0\in\Pi$. The map
$\ct:\Oz\to\Z$ takes the coefficient at exponent $0$.
\end{definition}

This is the standard ring, not the ring generated by the ordinals.
For example, $\sqrt2\,\omega$ is omnific, but $\omega+\pi$ is not.
The direct sum in \eqref{eq:Oz} is additive, not a direct product of
unital rings. Positive exponents stay positive under addition, so
$\Pi$ is an ideal and $\ct$ is a ring retraction.

\begin{proposition}[Integer part]\label{prop:floor}
The ring $\Oz$ is discretely ordered with least positive element $1$.
For every $x\in\No$ there is exactly one $z\in\Oz$ such that
$z\leq x<z+1$. This $z=\lfloor x\rfloor_{\Oz}$ is parameter-free
definable from $x$ in set theory.
\end{proposition}
\begin{proof}
Write, by splitting the normal form,
$x=P+r+\varepsilon$, where $P\in\Pi$, $r\in\R$, and
$\varepsilon$ is infinitesimal. Put
\begin{equation}\label{eq:floor}
 \lfloor x\rfloor_{\Oz}=P+
 \begin{cases}
 r-1,&r\in\Z\ \text{and}\ \varepsilon<0,\\
 \lfloor r\rfloor,&\text{otherwise}.
 \end{cases}
\end{equation}
A nonzero positive element of $\Pi$ is larger than every real, since
its leading exponent is positive. Hence every positive element of
$\Oz$ is at least $1$. Formula~\eqref{eq:floor} directly verifies the
two inequalities. If two omnific integers had the required property,
their nonzero difference would have absolute value less than $1$,
contradicting discreteness. Both the normal-form formula and the unique
integer-part property are set-theoretic definitions.
\end{proof}

\begin{remark}
The correction when $r$ is integral and $\varepsilon<0$ is essential:
$\lfloor\omega-\omega^{-1}\rfloor_{\Oz}=\omega-1$, not $\omega$.
The proposition is classical \cite{Conway,LM}; its proof is included
to make subsequent definability closure entirely explicit.
\end{remark}

\section{Which notion of definability?}
\label{sec:definitions}

\subsection{Ambient definability with parameters}

\begin{definition}\label{def:ambient}
Let $M\models\ZFC$ be a set model, and let $A$ be an external collection
of allowed elements of $M$. An element $a\in M$ is \emph{ambiently
$A$-definable in $M$} if some standard first-order formula
$\varphi(v,\bar u)$ and some finite tuple $\bar p$ from $A$ satisfy
\[
 M\models\forall v\,[\varphi(v,\bar p)\leftrightarrow v=a].
\]
Write $D_A^M=\dcl_{(M,\in)}(A)$, and define
\[
 K_A^M=D_A^M\cap\No^M,\qquad I_A^M=D_A^M\cap\Oz^M.
\]
When $A=\varnothing$, say \emph{parameter-free definable}. Superscripts
are omitted when the universe is fixed.
\end{definition}

For a set model, $D_A^M$ is an external set, definable using the external
satisfaction relation of $M$. In particular,
\[
 |D_A^M|\leq\max(\aleph_0,|A|).
\]
This estimate is external. It does not put into $M$ a list of all the
elements definable in $M$.

For the intended universe $V$, the same terminology is a metatheoretic
scheme: for each particular formula we can ask whether it uniquely
defines a given object. We do not add a first-order truth predicate for
$V$, and do not assume that the full collection $D_\varnothing^V$ is
an internally available set or class.

\subsection{Ordinal parameters}

\begin{definition}\label{def:OD}
A surreal or omnific integer is \emph{ordinal definable} if its canonical
set code is in $\OD$. Define
\[
 \KH=\No\cap\OD,\qquad \IH=\Oz\cap\OD.
\]
The standard internal definition of $\OD$ says that an object is uniquely
definable over some set structure $(V_\theta,\in)$ using finitely many
ordinal parameters below $\theta$. Reflection identifies this with the
usual notion of definability in $V$ from ordinal parameters. The extra
ordinal $\theta$ is allowed; no truth predicate for $V$ is required.
\end{definition}

Both $\OD$ and $\HOD$ are definable classes of ZFC, unlike a uniform
truth-based relation between arbitrary formulas and their unique
solutions in $V$. HOD is the transitive class of sets whose transitive
closures, including the sets themselves, consist of ordinal-definable
sets; it is an inner model of ZFC with a canonical definable well-order
\cite{Jech,GoldbergPoveda}. Every ordinal is ordinal definable, using
itself as an allowed parameter. This alone shows why $\KH$ need not
be a countable collection.

\subsection{Intrinsic ordered-field definability}

\begin{theorem}[Classical language contrast]\label{thm:purefield}
In the pure ordered-field language, the parameter-free definable
surreals are exactly the real algebraic numbers. Their intersection
with $\Oz$ is exactly $\Z$. More generally, in a real closed surreal
field $F$,
\[
 \dcl_F(A)=\text{the relative real closure of }\Q(A)\text{ in }F.
\]
\end{theorem}
\begin{proof}
Quantifier elimination for real closed fields reduces a definable
singleton to a semialgebraic singleton over $\Q(A)$ \cite{Marker}.
If a point is not a zero of any of the finitely many nonzero polynomials
appearing in a quantifier-free defining formula, their signs are locally
constant, so the formula cannot isolate that point. Thus the point is
algebraic over $\Q(A)$. Conversely, an algebraic element is specified as
the $k$th real root of its polynomial, for a fixed finite $k$.
Real algebraic elements of $\No$ lie in its standard real subfield.
A real is omnific exactly when it is an integer.
\end{proof}

For the full proper-class field, the argument is a formula-by-formula
application of the same quantifier-elimination theorem. In particular,
$\omega$ is \emph{ambiently} parameter-free definable, but is not
parameter-free definable in the pure field. Nor does adding the word
``omnific'' to our prose add a predicate to that language. The structure
$(\No;+,-,\cdot,<,0,1,\Oz)$ is a different structure, whose definable
closure is not classified here.

\subsection{Provable descriptions and computability}

For a recursively axiomatized theory $T$, one can restrict to formulas
$\varphi$ for which $T\vdash\exists!x\,(\No(x)\wedge\varphi(x))$.
These are \emph{$T$-provably unique descriptions}. Their interpretation
in a model of $T$ need not be invariant between different models.
The statement ``$T$ proves uniqueness'' is not a replacement for choosing
the semantic universe.

A genuine effective presentation of a surreal sign sequence gives an
ambient definition when the presentation uniquely determines a
well-founded sequence. The converse fails already for real numbers:
the characteristic real of a noncomputable but arithmetically definable
set is definable. We therefore do not identify definability with any
one of the computable-surreal representations catalogued by the
repository, or with a numerical approximation algorithm.

\section{The field and integer part of ambiently definable numbers}
\label{sec:hulls}

In this section the superscript $M$ refers to a fixed transitive model
when external semantics are used. All operations on its surreals and
normal forms are its internal operations.

\begin{lemma}[Unique-output closure]\label{lem:closure}
Suppose a set-theoretic formula defines a unique output $F(\bar x)$
from each input in its domain. If every component of $\bar a$ lies
in $D_A$, then $F(\bar a)\in D_A$.
\end{lemma}
\begin{proof}
Choose definitions for the finitely many components of $\bar a$.
Existentially quantify those components, assert their defining formulas,
and assert the graph formula for $F$. This is a single formula using
only finitely many parameters from $A$, with a unique solution.
\end{proof}

\begin{theorem}[Definable field and integer part]\label{thm:hull}
The ambiently $A$-definable surreals $K_A$ form a real closed field.
The ring $I_A=K_A\cap\Oz$ is a discretely ordered integer part of $K_A$,
and
\begin{equation}\label{eq:frac}
 K_A=\Frac(I_A).
\end{equation}
The field $K_A$ is an elementary substructure of the ambient surreal
ordered field, in the pure ordered-field language.
\end{theorem}
\begin{proof}
The constants $0,1$, addition, subtraction, multiplication, and inverse
away from zero are uniquely definable set-theoretic operations.
Lemma~\ref{lem:closure} proves the field axioms by restriction.
Every positive element has a unique positive square root in $\No$,
so that square root belongs to $K_A$. For an odd-degree polynomial
with coefficients in $K_A$, real closedness of $\No$ gives at least
one real root; its finite set of real roots has a least member, uniquely
definable from the coefficient tuple. That root lies in $K_A$.
The square-root and odd-degree-root criterion proves real closedness.
This is an application of the classical real-closed-field theorem,
not a new proof of real closedness of $\No$.

The floor from Proposition~\ref{prop:floor} belongs to $D_A$ whenever
its input does. Hence every $x\in K_A$ has its omnific floor in $I_A$.
Discreteness and the subring property restrict from $\Oz$.

For the fraction-field assertion, let $0\ne x\in K_A$, and put
$S=\supp(x)$. The whole normal form and $S$ are definable from $x$.
There exists a positive ordinal $\delta$ such that
\[
 \delta+g>0\qquad(g\in S)
\]
by Lemma~\ref{lem:ordinalbound}. Choose the \emph{least} such positive
ordinal. It is definable from $x$, so $\delta\in K_A$. Thus
$q=\omega^\delta$ and $p=qx$ belong to $K_A$. Their supports are
strictly positive, so $p,q\in I_A$, and $x=p/q$. The zero case is
$0=0/1$. The reverse inclusion follows because $K_A$ is a field.
Finally, every inclusion of real closed fields is elementary by
quantifier elimination \cite{Marker}.
\end{proof}

\begin{remark}
The proof does not suppose that each exponent of an $A$-definable
normal form is $A$-definable. It only uses definability of the
\emph{entire support} to choose one ordinal bound. Confusing these
two assertions would invalidate the fraction-field argument.
The equality $\Frac(\Oz)=\No$ itself is not claimed as new; the point
here is that denominator clearing can be chosen uniformly enough to
preserve ambient definability.
\end{remark}

\begin{corollary}[Open induction]\label{cor:openinduction}
The nonnegative part of $I_A$ satisfies induction for every
quantifier-free formula in the ordered-ring language, with parameters
from $I_A$.
\end{corollary}
\begin{proof}
A quantifier-free formula with parameters in $I_A$ defines a
semialgebraic subset of the real closed field $K_A$, hence a finite
union of intervals and points. Every nonempty intersection of such a
set with $(I_A)_{\geq0}$ has a least member. Indeed, the integer-part
property supplies the first integer at or immediately above each
finite left endpoint, with a one-unit adjustment for a strict endpoint;
an interval unbounded to the left is treated by the lower bound $0$.
Take the least of the finitely many nonempty pieces.
If an induction instance failed, the set of counterexamples would
therefore have a least member. It is not $0$, and its predecessor is
nonnegative and satisfies the formula, contradicting the inductive
step. This is the standard integer-part route to open induction
\cite{Shepherdson,LM}.
\end{proof}

\begin{proposition}[Finite ambiguity does not enlarge the hull]
\label{prop:algebraicity}
An element of $\No$ belonging to a finite ambiently $A$-definable set
of surreals is itself ambiently $A$-definable. In particular, ambient
model-theoretic algebraicity and definability coincide on the surreal
sort. Every surreal algebraic over $K_A$ already belongs to $K_A$.
\end{proposition}
\begin{proof}
Order the finite candidate set by the surreal order. Each member is
its $k$th member for a fixed finite $k$; this property is expressible
by a formula with the same parameters. For the final assertion, the
roots of a nonzero polynomial over $K_A$ form such a finite definable
set. Equivalently, a real closed field has no proper algebraic ordered
extension.
\end{proof}

This finite-ordering argument is a particularly transparent case of
model-theoretic algebraicity; broader set-theoretic questions about
algebraicity and implicit definability are studied in
Hamkins--Leahy \cite{HamkinsLeahy}. It concerns a genuinely finite
candidate set, not an internally nonstandard finite set viewed from
outside a nonstandard model.

\begin{proposition}[Uniform constructions]\label{prop:uniform}
The field $K_A$ is closed under the Conway monomial map and under
forming the simplest separator of an $A$-definable set-sized separated
cut. It is also closed under the sum of a valid normal form when the
\emph{whole indexed normal form} is $A$-definable.
\end{proposition}
\begin{proof}
Each construction has a unique set-theoretically specified output.
Apply Lemma~\ref{lem:closure}. For a cut, both option sets must be
jointly available as $A$-definable sets, not merely have individually
$A$-definable elements.
\end{proof}

\section{The ordinal-definable surreal field is the HOD field}
\label{sec:HOD}

The canonical sign representation makes hereditary definability much
more transparent than it is for arbitrary sets.

\begin{theorem}[HOD identification]\label{thm:HOD}
With canonical sign-sequence representatives,
\begin{equation}\label{eq:HODid}
 \KH=\No\cap\OD=\No\cap\HOD=\No^{\HOD},
 \qquad
 \IH=\Oz\cap\HOD.
\end{equation}
The field $\KH$ is real closed, contains every ordinal, and is
closed under the Conway monomial map. The ring $\IH$ is an integer
part with $\Frac(\IH)=\KH$.
\end{theorem}
\begin{proof}
Let $s:\alpha\to\{-,+\}$ be ordinal definable. Every element in its
transitive closure is built from ordinals and the two fixed sign
symbols by the fixed ordered-pair construction. All such elements are
ordinal definable. Since $s$ itself is ordinal definable, $s\in\HOD$.
The converse is part of the definition of HOD.

HOD contains every ordinal and is transitive. Its surreal sign
sequences are consequently exactly the ambient sign sequences that
belong to HOD. This gives the equality of carriers in
\eqref{eq:HODid}. Their ordered-field structures agree: comparison is
computed from the same signs; canonical options are the same initial
segments; and the recursive arithmetic on these set-sized option
closures is absolute between a transitive inner model and the ambient
universe. For the monomial map one may use its recursive cut with
integer and dyadic multipliers, so no newly added reals enter the
recursion. Thus the carrier identification is an ordered-field
identification, not just a coincidence of underlying sets.

Alternatively, all closure assertions follow directly by substituting
ordinal definitions into the unique-output constructions in
Theorem~\ref{thm:hull}. The ordinals used by those substitutions are
still permitted parameters. Proposition~\ref{prop:floor} and the same
least-ordinal denominator clearing prove the claims about $\IH$.
\end{proof}

\begin{lemma}[Support of an ordinal-definable number]
\label{lem:HODsupport}
If $x\in\KH$ has normal form \eqref{eq:normal}, then for every
$\xi<\lambda$,
\[
 r_\xi\in\R\cap\HOD,\qquad g_\xi\in\KH.
\]
In fact, $\NF(x)\in\HOD$, and
\begin{equation}\label{eq:NFcriterion}
 x\in\KH\quad\Longleftrightarrow\quad\NF(x)\in\HOD.
\end{equation}
\end{lemma}
\begin{proof}
The entire normal form is uniquely definable from $x$, so is ordinal
definable. An individual entry is uniquely definable from $x$ and the
ordinal index $\xi$. Its real coefficient is ordinal definable, hence
hereditarily ordinal definable under the usual real coding. Its
surreal exponent belongs to HOD by Theorem~\ref{thm:HOD}. The length,
indices, and the fixed tuple coding are also hereditarily ordinal
definable. Hence the whole sequence belongs to HOD. Conversely,
$\NF(x)\in\HOD$ implies that it is ordinal definable, so its unique
surreal value is ordinal definable.
\end{proof}

The last equivalence requires the whole sequence to lie in HOD. The
strictly weaker condition ``every coefficient and exponent lies in
HOD'' is insufficient; Section~\ref{sec:uniformity} supplies explicit
counterexamples, even with a HOD support and rational coefficients.
In particular, $\KH$ is not being identified with the unrestricted
ambient Hahn field containing \emph{all} externally assembled series
with coefficients and exponents in HOD.

\begin{corollary}[No algebraic defect]\label{cor:noalgebraic}
Every $x\in\No\setminus\KH$ is transcendental over $\KH$. The
inclusion $\KH\subseteq\No$ is elementary in the ordered-field
language, whether or not $V=\HOD$.
\end{corollary}
\begin{proof}
Apply Proposition~\ref{prop:algebraicity} with ordinal parameters,
and then quantifier elimination.
\end{proof}

\begin{proposition}[The HOD field is not intrinsically definable when proper]
If $\KH\ne\No$, no ordered-field formula, even with finitely many
surreal parameters, defines $\KH$ inside $\No$.
\end{proposition}
\begin{proof}
An infinite definable subset of a real closed field is a finite union
of points and intervals, so an infinite definable subfield contains
an interval. Taking differences, it contains $(0,c)$ for some
$c>0$ belonging to the subfield. For any $x>0$, the number
$y=1/(x+2/c)$ lies in $(0,c)$, hence in the subfield. Then
$x=1/y-2/c$ belongs to it as well. Thus every infinite definable
subfield is the entire field.
\end{proof}

\section{A reversible code in a fixed omnific interval}
\label{sec:code}

For each ordinal $\beta$, viewed as a surreal ordinal, define
\begin{equation}\label{eq:eps}
 e_\beta=\omega^{-(\beta+1)}.
\end{equation}
The $e_\beta$ form a strictly decreasing class of positive infinitesimal
surreals. In particular $0<e_\beta\leq\omega^{-1}<1/2$.
Every $e_\beta$ is ordinal definable. For a sign sequence
$s:\alpha\to\{-,+\}$ put
\[
 c_s(\beta)=\begin{cases}1,&s(\beta)=-,\\2,&s(\beta)=+.\end{cases}
\]

\begin{definition}[Self-delimiting omnific code]\label{def:code}
Define
\begin{equation}\label{eq:code}
 \Code(s)=\omega+
       \sum_{\beta<\alpha}c_s(\beta)\omega^{e_\beta}
       +3\omega^{e_\alpha},\qquad\alpha=\bd(s).
\end{equation}
The coefficient $3$ is an end marker, not a third sign value.
\end{definition}

\begin{theorem}[Reversible fixed-leading-term coding]\label{thm:code}
Formula~\eqref{eq:code} defines a parameter-free set-theoretic injection
$\Code:\No\to\Cell$. Its image is set-theoretically definable, and
there is a parameter-free definable inverse $\Decode$ on that image.
For every model $M$ and every allowed parameter family $A$,
\begin{equation}\label{eq:degree}
 s\in D_A^M\quad\Longleftrightarrow\quad\Code^M(s)\in D_A^M.
\end{equation}
In particular, $s$ and $\Code(s)$ are mutually definable and have the
same ordinal-definability status.
\end{theorem}
\begin{proof}
The exponents in \eqref{eq:code} occur in strictly decreasing order:
first $1$, then $e_\beta$ for $\beta\leq\alpha$. They form a set,
because $\alpha$ is an ordinal. Each coefficient is nonzero and real,
so the expression is a valid Conway normal form. Every exponent is
positive. Thus the value is an omnific integer with constant term zero.
Its leading term is exactly $\omega$.

The tail $T=\Code(s)-\omega$ is positive. Its leading exponent is
$e_0=\omega^{-1}$, including the empty-sign case, in which the leading
tail coefficient is $3$. Since $e_0<1/2$, comparison of normal forms
gives $T<\omega^{1/2}$. Hence $\Code(s)\in\Cell$.

To decode, take the unique normal form of $z$. Verify that its leading
term is $\omega$, that its remaining exponents are exactly
$e_\beta$ for an initial ordinal segment $\beta\leq\alpha$, that
the last coefficient is $3$, and that every preceding coefficient is
$1$ or $2$. The map $\beta\mapsto e_\beta$ is injective, so the end
marker uniquely determines $\alpha$. At each $\beta<\alpha$, recover
$-$ from coefficient $1$ and $+$ from coefficient $2$. This produces
exactly one sign sequence $s$ and satisfies $\Decode(\Code(s))=s$.

These are set-theoretic definitions: the assertion that there exists
an ordinal $\alpha$ and the specified indexed normal form is a
first-order assertion about sets. It neither quantifies over class
sequences nor needs a field-intrinsic way to recognize $\omega$.
The unique-output closure lemma applied to both maps proves
\eqref{eq:degree}. The same proof is a theorem internal to any model;
external definability uses standard formulas for the two maps.
\end{proof}

\begin{example}
The first three codes are
\[
 \Code(0)=\omega+3\omega^{\omega^{-1}},
\]
\[
 \Code(1)=\omega+2\omega^{\omega^{-1}}+3\omega^{\omega^{-2}},
 \qquad
 \Code(-1)=\omega+\omega^{\omega^{-1}}+3\omega^{\omega^{-2}}.
\]
The code for $0$ is not $0$. The coding theorem makes no assertion of
arithmetic compatibility.
\end{example}

\begin{corollary}[Congruence and leading-data blindness]\label{cor:blind}
Every code has leading exponent $1$, leading coefficient $1$, constant
term $0$, and coefficients from $\{1,2,3\}$. Moreover, for every
ordinary positive integer $n$,
\[
 \Code(s)\in n\Oz.
\]
These invariants therefore do not bound ambient definability complexity.
\end{corollary}
\begin{proof}
Division by an ordinary $n>0$ scales the real coefficients and leaves
all exponents positive, so the quotient is still omnific. The remaining
assertions follow from the normal form. Since every surreal is
recoverable from its code, the listed common invariants cannot
separate definability degrees.
\end{proof}

\begin{proposition}[A complementary two-term code]\label{prop:twotermcode}
There is also a parameter-free definable order isomorphism
\[
 \mathsf T:\No\longrightarrow
 \Cell_2:=\{\omega+\omega^g:0<g<1/2\}\subseteq\Cell.
\]
Every value has exactly two nonzero normal-form terms, both with
coefficient $1$, and is mutually ambiently definable with its input.
\end{proposition}
\begin{proof}
The rationally defined map
\[
 h(x)=\frac14\left(1+\frac{x}{1+|x|}\right)
\]
is an increasing bijection from $\No$ onto $(0,1/2)$. For its
inverse, put $u=4g-1$ and use $x=u/(1-|u|)$. Therefore
$\mathsf T(x)=\omega+\omega^{h(x)}$ has the stated normal form,
is strictly increasing, and has exactly the stated image. Extracting
the lower exponent from the unique normal form and applying the
inverse of $h$ recovers $x$ by a parameter-free set-theoretic
operation.
\end{proof}

The two codes have different purposes. The elementary two-term
code allows an arbitrary surreal exponent to carry all the input
information. The sign code $\Code$ instead fixes its exponents in
an ordinal-definable template, so that all information is stored in
the assignment of a finite coefficient alphabet. It is the latter
feature that will prove the failure of termwise HOD closure.

\begin{remark}[An interval is still a proper class]
Although $\Cell$ is bounded by two explicitly given surreals, it is a
proper class in the ambient universe. Otherwise its subset
$\ran(\Code)$ would be a set and Replacement applied to $\Decode$
would make all of $\No$ a set. Ordered boundedness is not a size bound
for the full surreal field.
\end{remark}

\section{One interval detects the entire set-theoretic universe}
\label{sec:global}

We first spell out the standard set-coding step needed for the global
consequences. This step is not an original bi-interpretation claim.

\begin{lemma}[Every set is recoverable from some surreal]
\label{lem:setcoding}
There is a parameter-free partial set-theoretic function $E$ on
surreal sign sequences such that every set is $E(s)$ for some $s$.
No parameter-free selection of such an $s$ from the set being coded
is asserted.
\end{lemma}
\begin{proof}
For a set $a$, Choice gives a bijection $f:\kappa\to\TC(\{a\})$
for some ordinal $\kappa$. Transport membership to a well-founded
extensional relation $R$ on $\kappa$, and record the unique index
$\rho$ with $f(\rho)=a$.

Fix a parameter-free class pairing of ordinals
$p:\Ord^2\to\Ord$. One construction well-orders ordinal pairs first
by their maximum and then lexicographically, and uses the order type
of each predecessor set. Every predecessor collection is a set, so
this defines a class bijection. Form the set of ordinals
\[
 B=\{p(0,\kappa),p(1,\rho)\}
       \cup\{p(2,p(\xi,\eta)):(\xi,\eta)\in R\}.
\]
Let $\lambda=\sup\{\gamma+1:\gamma\in B\}$ and take the sign
sequence of length $\lambda$ whose positive positions are exactly
$B$. From this sequence recover $B$, then $\kappa,\rho,R$, and finally
the value at $\rho$ of the unique Mostowski collapse. These operations
are parameter-free and unique whenever the code is valid. For invalid
codes leave $E$ undefined. Mostowski collapse is standard
\cite{Jech}; the same pointed-graph mechanism is present in the
announced surreal/set-theory dictionary \cite{CHY}.
\end{proof}

\begin{theorem}[The ordinal-definability test]\label{thm:globalOD}
The following statements are equivalent in ZFC:
\begin{enumerate}[label=\textup{(\roman*)}]
\item $V=\HOD$;
\item every surreal is ordinal definable;
\item every omnific integer is ordinal definable;
\item every member of $\Cell$ is ordinal definable;
\item every value of the restricted code $\Code$ is ordinal definable.
\end{enumerate}
\end{theorem}
\begin{proof}
The forward implications through (v) are immediate from the inclusions.
If (v) holds, recover each $s$ from $\Code(s)$ by
Theorem~\ref{thm:code}; hence every surreal is ordinal definable.
For any set $a$, choose a surreal $s$ with $E(s)=a$ by
Lemma~\ref{lem:setcoding}. Its canonical sign code lies in HOD by
Theorem~\ref{thm:HOD}. More explicitly, the associated ordinal set
$B$ lies in HOD, and HOD decodes the same well-founded extensional
relation and its Mostowski collapse. Thus $a\in\HOD$.
Since $a$ was arbitrary, $V=\HOD$.
\end{proof}

\begin{corollary}[Two terms already detect HOD]
The equivalent conditions in Theorem~\ref{thm:globalOD} remain
equivalent if condition \textup{(iv)} is replaced by ``every member
of $\Cell_2$ is ordinal definable.''
\end{corollary}
\begin{proof}
Use the two mutually definable maps in
Proposition~\ref{prop:twotermcode} to recover every surreal from
such a two-term omnific integer, and then apply
Theorem~\ref{thm:globalOD}.
\end{proof}

\begin{remark}
There is no Choice-based definable selection hidden in the converse.
We choose a code only to prove its existence; hypothesis (ii) says
\emph{every} such surreal code is ordinal definable. This is much
weaker as a proof requirement than constructing a canonical code from
an arbitrary set in a universe not satisfying $V=\HOD$.
\end{remark}

\begin{theorem}[The pointwise-definability test]\label{thm:pointwise}
Let $M\models\ZFC$ be any set model. In external semantics, the
following are equivalent:
\begin{enumerate}[label=\textup{(\roman*)}]
\item $M$ is pointwise definable;
\item every member of $\No^M$ is parameter-free definable in $M$;
\item every member of $\Oz^M$ is parameter-free definable in $M$;
\item every member of $\Cell^M$ is parameter-free definable in $M$.
\end{enumerate}
\end{theorem}
\begin{proof}
Only (iv)$\Rightarrow$(i) needs proof. Each $s\in\No^M$ is recovered
by a fixed standard formula from $\Code^M(s)\in\Cell^M$, so (iv)
implies (ii). Internally to $M$, every set has a code as in
Lemma~\ref{lem:setcoding}. For any fixed $a\in M$, choose such a
code $s\in M$. A standard formula defines $s$ by (ii), and composing
it with the standard decoder formula defines $a$. This does not
require $M$'s internal relation to be externally well-founded:
the collapse operation is the one that $M$ proves exists uniquely.
\end{proof}

Pointwise definable models of ZFC exist whenever ZFC is consistent;
Hamkins--Linetsky--Reitz also prove strong extension theorems
\cite{HLR}. Their existence theorem is prior work. The new feature of
Theorem~\ref{thm:pointwise} is the restriction of the test to a single
fixed-leading-term omnific interval. It also explains why an external
counting argument cannot be converted into an internal proof that
some member of $\Oz^M$ is undefinable. The formulas that define
particular elements are not presented to $M$ as a surjective internal
list.

\section{From one missing number to proper-class transcendence}
\label{sec:transcendence}

This section contains the strongest algebraic result. We first isolate a
support argument that does not involve definability at all.

\begin{lemma}[Disjoint support cosets]\label{lem:cosets}
Let $K\subseteq\No$ be a subfield, and let $G\subseteq(\No,+)$ be an
additive subgroup such that
\[
 \supp(a)\subseteq G\qquad(0\ne a\in K).
\]
Suppose $(g_i)_{i\in J}$ is a family such that every nonzero finitely
supported integer combination satisfies
\begin{equation}\label{eq:modind}
 \sum_i m_i g_i\notin G.
\end{equation}
Then the family $(\omega^{g_i})_{i\in J}$ is algebraically independent
over $K$. The assertion applies to a class family when the condition
is read on each finite subfamily.
\end{lemma}
\begin{proof}
Take a nonzero polynomial in finitely many distinct variables,
\[
 P(Y_1,\ldots,Y_n)=\sum_{\nu\in F}a_\nu Y_1^{\nu_1}\cdots
 Y_n^{\nu_n},
\]
where $F\subseteq\N^n$ is finite and the displayed coefficients are
nonzero. Since $\omega^u\omega^v=\omega^{u+v}$, the normal-form
support of the term at multi-index $\nu$ after substitution is
\begin{equation}\label{eq:supportshift}
 \supp(a_\nu)+\sum_{j=1}^n\nu_jg_j
 \ \subseteq\ G+\sum_{j=1}^n\nu_jg_j.
\end{equation}
For $\nu\ne\mu$, these cosets are disjoint. Otherwise subtracting an
exponent in their intersection would give
$\sum_j(\nu_j-\mu_j)g_j\in G$, contrary to \eqref{eq:modind}.
Every displayed term is nonzero, and its coefficients cannot be
cancelled by any other term. The finite sum is therefore nonzero.
This proves the required absence of polynomial relations. No infinite
polynomial, class sum, or interchange of summations is used.
\end{proof}

\begin{theorem}[Support-subfield amplification]\label{thm:amplify}
Let $K$ be a proper subfield of $\No$ satisfying
\begin{equation}\label{eq:Khyp}
 \omega\in K,\qquad e_\alpha\in K\ (\alpha\in\Ord),\qquad
 \supp(a)\subseteq K\ (0\ne a\in K),
\end{equation}
where $e_\alpha=\omega^{-(\alpha+1)}$.
For any $t\in(0,1)\setminus K$, the class family
\begin{equation}\label{eq:indfam}
 y_\alpha=\omega^{t e_\alpha},\qquad
 z_\alpha=\omega+y_\alpha\qquad(\alpha\in\Ord)
\end{equation}
is algebraically independent over $K$, both for the $y$ family and
for the $z$ family. Each $z_\alpha$ lies in $\Cell$ and has exactly
two nonzero normal-form terms, both with coefficient $1$.
Such a $t$ exists for every proper $K$ satisfying \eqref{eq:Khyp}.
\end{theorem}
\begin{proof}
First, the family $(e_\alpha)$ is linearly independent over $\Q$.
Indeed, a finite rational linear combination is already a normal
form with distinct exponents $-(\alpha+1)$; after deleting zero
coefficients it cannot vanish. Hence for distinct
$\alpha_1,\ldots,\alpha_n$ and integers $m_j$ not all zero,
\[
 c=\sum_{j=1}^n m_j e_{\alpha_j}\in K\setminus\{0\}.
\]
If $\sum_jm_jte_{\alpha_j}=tc$ belonged to $K$, division by $c$
would give $t\in K$, a contradiction. Thus the exponents
$te_\alpha$ satisfy \eqref{eq:modind} for the additive subgroup
$G=(K,+)$. Lemma~\ref{lem:cosets} proves algebraic independence of
the $y_\alpha$.

Translation by $\omega\in K$ is a polynomial automorphism on each
finite tuple of variables. Therefore a nonzero relation among the
$z_\alpha=\omega+y_\alpha$ would give a nonzero relation among the
$y_\alpha$, which is impossible.

Finally,
\[
 0<te_\alpha<e_\alpha\leq\omega^{-1}<1/2.
\]
Consequently $\omega^{te_\alpha}$ is a positive, purely infinite
omnific monomial strictly below $\omega^{1/2}$. The exponents $1$
and $te_\alpha$ are distinct and positive, so the asserted two-term
normal form is exact and $z_\alpha\in\Cell$.

For existence, take $x\in\No\setminus K$. Then $u=|x|>0$ is not
in $K$, since otherwise $x=u$ or $x=-u$ would be in $K$. Set
$t=u/(1+u)$. It lies strictly between $0$ and $1$. If $t\in K$,
then $u=t/(1-t)\in K$, again a contradiction.
\end{proof}

\begin{theorem}[The HOD-surreal transcendence dichotomy]
\label{thm:dichotomy}
Exactly one of the following occurs:
\begin{enumerate}[label=\textup{(\roman*)}]
\item $V=\HOD$ and $\No=\KH$;
\item $V\ne\HOD$, and for any $t\in(0,1)\setminus\KH$, the family
\[
 \left(\omega+\omega^{\,t\omega^{-(\alpha+1)}}\right)_{\alpha\in\Ord}
\]
is algebraically independent over $\KH$ and is contained in $\Cell$.
\end{enumerate}
In case \textup{(ii)}, for every cardinal $\kappa$ there is an
algebraically independent subset of $\Cell$ of cardinality $\kappa$
over $\KH$. In particular, the transcendence defect has no set-sized
cardinal bound.
\end{theorem}
\begin{proof}
Theorem~\ref{thm:globalOD} identifies $V\ne\HOD$ with
$\KH\ne\No$. The field $\KH$ contains $\omega$ and every
$e_\alpha$ because these are ordinal definable. Its support property
is Lemma~\ref{lem:HODsupport}. Apply Theorem~\ref{thm:amplify}.
Restricting the class family to $\alpha<\kappa$ gives a set family
by Replacement. Algebraic independence implies distinctness, hence
that set has cardinality $\kappa$.
\end{proof}

\begin{remark}[What ``proper-class transcendence'' means]
The theorem gives a concrete class injection with finite polynomial
independence, definable from one parameter $t$. It does not choose a
transcendence basis for $\No$, invoke Global Choice, or assign an
ordinary cardinal to a proper class. Each individual $z_\alpha$ is
transcendental over $\KH$, but this pointwise assertion alone would
be much weaker than the simultaneous algebraic independence proved
above. In a set model, the assertions have their internal cardinal
and ordinal meaning, not an assertion that its externally set-sized
carrier is literally a proper class in the metatheory.
\end{remark}

\begin{corollary}[Localization at other definable scales]
\label{cor:localize}
Assume $V\ne\HOD$. For every $p\in\IH$, every positive
$\gamma\in\KH$, and every $t\in(0,1)\setminus\KH$, the family
\begin{equation}\label{eq:localfam}
 w_\alpha=p+\omega^{\gamma t e_\alpha}\qquad(\alpha\in\Ord)
\end{equation}
is algebraically independent over $\KH$ and lies in
$(p,p+\omega^\gamma)\cap\Oz$.
\end{corollary}
\begin{proof}
For every nonzero finite integer combination $c$ of the $e_\alpha$,
we have $c\in\KH\setminus\{0\}$. If $\gamma tc\in\KH$, division
by $\gamma c$ would put $t$ in $\KH$. The coset proof therefore
applies to the exponents $\gamma te_\alpha$. Translation by
$p\in\KH$ preserves independence. The exponents are positive and
less than $\gamma$, so the added monomials are omnific and lie
strictly between $0$ and $\omega^\gamma$.
\end{proof}

The widths in this corollary are infinite monomial widths. It would
be false to replace them indiscriminately by positive real widths:
an interval of width less than $1$ contains at most one omnific
integer. The result is localization at an arbitrary definable
\emph{infinite scale}, not density of omnific integers in $\No$.

\section{Uniform definability is not termwise definability}
\label{sec:uniformity}

There are two different pitfalls. Individually definable terms need
not form a definable sequence; conversely, a definable sequence need
not have parameter-free definable entries at arbitrary ordinal indices.

\begin{theorem}[Definable support and terms do not suffice]
\label{thm:termwise}
Assume $V\ne\HOD$. There is a valid normal form $q$ satisfying all
of the following:
\begin{enumerate}[label=\textup{(\roman*)}]
\item $q\in\Cell\setminus\IH$;
\item its support, as a whole set, belongs to HOD;
\item every coefficient is in $\{1,2,3\}$, and every individual
monomial term belongs to $\KH$;
\item the whole indexed coefficient/exponent sequence does not
belong to HOD.
\end{enumerate}
Thus $\KH$ is not closed under arbitrary ambient Hahn-summable
families of its elements, even when the support belongs to HOD.
\end{theorem}
\begin{proof}
Choose $s\notin\KH$ by Theorem~\ref{thm:globalOD}, and let
$\alpha=\bd(s)$ and $q=\Code(s)$. The reversible code gives
$q\notin\IH$. Its support is the set
\[
 \{1\}\cup\{e_\beta:\beta\leq\alpha\},
\]
which is definable from the ordinal $\alpha$ and whose members all
belong to HOD. Hence the support itself belongs to HOD. Every
individual coefficient is one of three fixed rationals, and every
individual exponent is ordinal definable, so every individual term
is ordinal definable. If the whole normal form belonged to HOD,
Lemma~\ref{lem:HODsupport} would give $q\in\KH$, a contradiction.
The summable family here is simply the family of its distinct
monomial terms, so no additional summability issue occurs.
\end{proof}

\begin{remark}
In this example even the \emph{set of coefficient values} is a
finite HOD set. What fails is the definability of which coefficient
is assigned to which exponent. Saying ``the coefficients are
rational'' erases precisely the information carrying the original
sign sequence.
\end{remark}

\begin{proposition}[A definable integer can encode nondefinable entries]
\label{prop:collector}
There is a parameter-free ambiently definable $q_*\in\Cell$ whose
normal-form coefficients canonically encode every member of
$\R\cap\HOD$. In any universe or fixed transitive model in which
some member of $\R\cap\HOD$ is not parameter-free definable, $q_*$
has a normal-form coefficient that is not parameter-free definable.
The last hypothesis is essential and is not asserted in every model.
\end{proposition}
\begin{proof}
The set $R_H=\R\cap\HOD$ is parameter-free definable. Restrict the
canonical set-like HOD well-order to it, and let
$(r_\xi)_{\xi<\lambda}$ be the unique increasing enumeration. The
entire enumeration and its ordinal order type $\lambda$ are
parameter-free definable. Define the injective real function
\[
 h(r)=\frac32+\frac{r}{2(1+|r|)}\in(1,2).
\]
Its inverse on $(1,2)$ is parameter-free definable: if
$u=2h(r)-3$, then $r=u/(1-|u|)$. Form
\begin{equation}\label{eq:collector}
 q_*=\omega+\sum_{\xi<\lambda}h(r_\xi)\omega^{e_\xi}.
\end{equation}
This is a valid normal form and is parameter-free definable from the
whole enumeration. The set $R_H$ is nonempty, and the leading tail
exponent is $e_0<1/2$, so $q_*\in\Cell$. Every $r\in R_H$ occurs
in exactly one coefficient, after applying $h$.
If that $r$ is not parameter-free definable, neither is $h(r)$,
since the displayed inverse would otherwise define $r$.
\end{proof}

This does not contradict unique-output closure. Extracting the
coefficient at position $\xi$ uses $\xi$ as an additional parameter;
a parameter-free definition of $q_*$ need not define every such
ordinal position. Ordinal definability allows that parameter, which
is why all coefficients of an ordinal-definable surreal do belong
to HOD.

\section{The first birthday of a non-ordinal-definable surreal}
\label{sec:firstbirthday}

When $V\ne\HOD$, there is a genuine definable first place where
ordinal-definable signs cease to exhaust the construction. It is
not the birthday of a canonically selected nondefinable number.

\begin{definition}
Assuming $V\ne\HOD$, define
\begin{equation}\label{eq:firstlambda}
 \lambda_H=\min\{\alpha\in\Ord:
                 \mathcal P(\alpha)\nsubseteq\HOD\}.
\end{equation}
Existence follows from the set-of-ordinals coding in
Lemma~\ref{lem:setcoding}. All terms in this definition are
first-order definable classes or set constructions.
\end{definition}

\begin{theorem}[Exact birthday threshold]\label{thm:firstbirthday}
Assume $V\ne\HOD$. Then
\begin{equation}\label{eq:firstbd}
 \lambda_H=\min\{\bd(s):s\in\No\setminus\KH\}.
\end{equation}
The ordinal $\lambda_H$ is an infinite cardinal in both HOD and $V$.
It is itself a parameter-free definable ordinal, and therefore a
parameter-free definable omnific integer.
\end{theorem}
\begin{proof}
For a sign sequence of length $\alpha$, its set of positive positions
is a subset of $\alpha$, and the two objects are mutually definable
from each other and $\alpha$. Since $\alpha\in\HOD$, one belongs
to HOD exactly when the other does. This proves \eqref{eq:firstbd}.
Every finite subset of an ordinal is ordinal definable, so the
threshold is infinite.

Suppose first that $\lambda_H$ is not a cardinal of HOD. A bijection
in HOD from some $\kappa<\lambda_H$ onto $\lambda_H$ would pull a
non-HOD subset of $\lambda_H$ back to a non-HOD subset of $\kappa$.
Indeed, if the pullback were in HOD, its image under that same
bijection would be in HOD. This contradicts minimality.

Next suppose that $V$ has a bijection from an infinite cardinal
$\kappa<\lambda_H$ onto $\lambda_H$. Transport the ordinal order
to a well-order $R$ of $\kappa$ having order type $\lambda_H$.
The ordinal $\kappa$, being a cardinal in $V$, is also a cardinal
in HOD. Choose in HOD a bijection between $\kappa$ and
$\kappa\times\kappa$. It codes $R$ as a subset of $\kappa$.
Every subset of $\kappa$ belongs to HOD by minimality, so $R$ itself
belongs to HOD. Its well-order collapse in HOD is absolute and gives
a HOD bijection between $\kappa$ and $\lambda_H$, contradicting
that $\lambda_H$ is a HOD cardinal. Thus it is also a cardinal in $V$.

Finally, \eqref{eq:firstlambda} is a parameter-free unique definition
of $\lambda_H$. Ordinal surreals are omnific, giving the last claim.
\end{proof}

\begin{corollary}[No definable witness selection]
The set of non-ordinal-definable signs of length $\lambda_H$ is a
nonempty parameter-free definable set, but none of its members is
ordinal definable. In particular, there is no ordinal-definable
selection of one of those signs, or of one of their omnific codes.
\end{corollary}
\begin{proof}
The set is defined using the definable predicate HOD and the
parameter-free ordinal $\lambda_H$. Its nonemptiness follows from
the definition of that ordinal. A selected ordinal-definable member
would contradict the defining condition. An ordinal-definable code
would yield such a member by the decoder.
\end{proof}

\subsection{A current set-theoretic restriction on the threshold}

The threshold is the first cardinal where $V$ and HOD disagree about
subsets. It therefore inherits restrictions already proved in set
theory; these are \emph{not} new surreal theorems in substance.
Benhamou, Cummings, Goldberg, Hayut, and Poveda prove that a singular
strong-limit cardinal of uncountable cofinality cannot be the first
such disagreement, and, from a measurable cardinal, obtain a model
whose first disagreement is $\aleph_\omega$
\cite{BCHGP2026}. Hence their first theorem gives:

\begin{corollary}[Translation of a cited preprint theorem]
If $\lambda_H$ is singular and strong limit, then
$\operatorname{cf}(\lambda_H)=\omega$.
\end{corollary}
\begin{proof}
An infinite singular cardinal has infinite cofinality. Uncountable
cofinality is excluded by the cited compactness theorem, so its
cofinality must be $\omega$.
\end{proof}

The preprint was submitted on August 25, 2026. It is important not to
replace the cardinal conclusion of Theorem~\ref{thm:firstbirthday}
by an unsupported regular-cardinal claim: the cited relative
consistency result explicitly permits a singular first threshold.
The older repository discussion of the first new birthday across
universes provides a closely related framework
\cite{RepoBirthday}; here the inner universe is specifically HOD.

\section{Complexity-bounded definability and escape ordinals}
\label{sec:complexity}

The unrestricted phrase ``all parameter-free definable numbers'' is
external. A fixed syntactic complexity bound, by contrast, admits an
internal treatment. This gives a precise version of the familiar
intuition that definitions of bounded complexity can be diagonalized
against without creating a paradox.

\subsection{The bounded hierarchy}

Fix a standard natural number $n\geq1$. Use a standard syntactic
coding of formulas in the L\'evy hierarchy, with padding so that
$\Sigma_n$ formulas are also admitted at every higher level. There
is a partial truth formula $\Sat_{\Sigma_n}(e,x)$ for this fixed
fragment. The existence of each fixed-fragment truth definition is
standard; it does not provide one truth formula for all values of
$n$ simultaneously \cite{Jech,HLR}.

\begin{definition}
Let $S_n$ be the collection of unique solutions of parameter-free
$\Sigma_n$ formulas, and put $B_n=S_n\cap\No$. More explicitly,
$x\in S_n$ means that some code $e$ of a formula with one free
variable satisfies
\[
 \Sat_{\Sigma_n}(e,x)\ \wedge\
 \forall y\bigl(\Sat_{\Sigma_n}(e,y)\longrightarrow y=x\bigr).
\]
Each formula is required to have a unique solution among all sets,
not just among the surreals.
\end{definition}

\begin{lemma}
For every fixed standard $n$, both $S_n$ and $B_n$ are
parameter-free definable sets of ZFC.
\end{lemma}
\begin{proof}
For a formula code having a unique solution, send the code to that
solution; send every other code to a fixed default set. The
fixed-fragment truth formula defines this total function on a subset
of $\omega$, so Replacement gives its range as a set. Separation
removes the default value when necessary and restricts to actual
unique solutions. A further Separation restricts to surreal codes.
All definitions use the fixed standard integer $n$ only through its
finite numeral, so there are no free parameters.
\end{proof}

\begin{theorem}[Definability-complexity escape]\label{thm:escape}
For every fixed standard $n\geq1$, let
\begin{equation}\label{eq:escape}
 \delta_n=\sup\{\bd(x)+1:x\in B_n\},
 \qquad d_n=\text{the all-plus surreal of length }\delta_n.
\end{equation}
Then $d_n$ is a parameter-free definable omnific integer, but
$d_n\notin B_n$. There is a standard $m>n$ with $d_n\in B_m$.
Consequently the hierarchy $(B_n)$ does not eventually stabilize.
\end{theorem}
\begin{proof}
Replacement gives the set of birthdays and the ordinal supremum in
\eqref{eq:escape}. The all-plus sequence is then uniquely defined,
so $d_n$ is parameter-free definable and is an ordinal surreal.
If $d_n\in B_n$, the definition of the supremum would imply
\[
 \bd(d_n)+1=\delta_n+1\leq\delta_n,
\]
which is impossible. On the other hand, the parameter-free formula
just given for $d_n$ is a particular finite formula. After conversion
to prenex form and adding dummy quantifiers, it belongs to some
$\Sigma_m$ fragment. Thus $d_n\in B_m$. By nesting and its failure
to lie in $B_n$, we must have $m>n$.
\end{proof}

\begin{remark}[No unproved complexity estimate]
The theorem does not assert $m=n+1$, a closed expression for
$\delta_n$, or a computable method that outputs the full transfinite
sign sequence of $d_n$. The syntactic complexity of the normal-form,
birthday, and sign-coding formulas depends on coding conventions.
Only existence of a finite higher level is used. Nor is the sequence
of escape ordinals asserted to be cofinal in all ordinals.
\end{remark}

\subsection{Why the separate definitions cannot be collected uniformly}

\begin{proposition}[No definable set contains its own full hull]
\label{prop:nouniformset}
There is no ambiently $A$-definable set $S\subseteq\No$ such that
$K_A\subseteq S$. In particular, no ambiently $A$-definable map
$\omega\to\No$ has range containing all of $K_A$.
\end{proposition}
\begin{proof}
If such an $S$ were given, choose the least ordinal strictly above
all its birthdays and form its all-plus surreal representative $d$.
The construction is uniquely definable from $S$, so $d\in K_A$.
But its birthday is strictly larger than that of any member of $S$,
so $d\notin S$, a contradiction. For a definable map on $\omega$,
Replacement gives an $A$-definable range set, reducing to the first
assertion.
\end{proof}

\begin{corollary}[Failure of a uniform bounded-truth collector]
\label{cor:nouniformBn}
In an $\omega$-standard model of ZFC there is no single
parameter-free formula that, for each $n\geq1$, uniquely defines
the set $B_n$ above, uniformly in $n$.
\end{corollary}
\begin{proof}
A uniform formula would, by Replacement, give a parameter-free set
$B=\bigcup_{n\geq1}B_n$. Every parameter-free definable surreal is
uniquely defined by a finite formula and hence belongs to some
$B_n$. Thus $B$ would contain $K_\varnothing$, contrary to
Proposition~\ref{prop:nouniformset}. The $\omega$-standardness
hypothesis ensures that internal natural indices are precisely the
standard finite levels discussed here.
\end{proof}

This is the exact point at which the false ``least undefinable
surreal'' argument tries to use more truth than it has. Separate
formulas defining each fixed $B_n$ do not constitute one internal
formula defining the sequence of all these sets. The argument also
does not claim that the membership class of definable elements can
never accidentally be definable: in a pointwise definable model it
is simply the entire universe. What fails there is an internal
surjective set-sized list of those elements, not the external fact
that every one of them has some definition.

\section{Forcing can change the theory without adding reals}
\label{sec:forcing}

We now give a concrete source of models exhibiting the dichotomy.
The forcing facts are standard, but proofs of the relevant
homogeneity and closure steps are included. The result connects
ambient definability of real numbers to genuinely transfinite
surreal information without conflating the two.

\begin{lemma}[Homogeneity over a constructible ground]
\label{lem:forcingHOD}
Let $W$ be a transitive model of ZFC satisfying $V=L$, and let
$G$ be generic for an ordinal-definable weakly homogeneous forcing
$\mathbb P\in W$. Then, in $W[G]$,
\[
 \HOD^{W[G]}=W.
\]
The assertion has the usual relative-model interpretation; it does
not assume that a generic filter over the actual universe is a set
in that universe.
\end{lemma}
\begin{proof}
First let $A\subseteq\theta$ be ordinal definable in $W[G]$.
Membership in $A$ is expressed there by one formula
$\psi(\beta,\bar\alpha)$ with a finite ordinal tuple. For each
$\beta<\theta$, weak homogeneity implies that the top condition
decides $\psi(\check\beta,\check{\bar\alpha})$.
Indeed, if one condition forced the sentence and another forced its
negation, an automorphism could make their images compatible, while
fixing all check names for ordinals. Compatible conditions cannot
force contradictory sentences. Thus
\[
 A=\{\beta<\theta:1_{\mathbb P}\Vdash
                 \psi(\check\beta,\check{\bar\alpha})\}\in W
\]
by the definability of the forcing relation.

Every HOD set is coded within HOD by a set of ordinals: use its
canonical HOD well-order on a transitive closure and the pointed
membership-relation code. Such a code is itself ordinal definable,
so the previous paragraph places it in $W$. The absolute Mostowski
collapse in $W$ then recovers the original set. This proves
$\HOD^{W[G]}\subseteq W$.
Conversely, forcing does not change the ordinals or the constructible
hierarchy, so $L^{W[G]}=W$. Every constructible set is ordinal
definable in $W[G]$ using its position in the canonical constructible
well-order, and so are all members of its transitive closure.
Therefore $W\subseteq\HOD^{W[G]}$.
These are the standard homogeneity/HOD arguments \cite{Jech}.
\end{proof}

\begin{theorem}[Prescribed regular first birthday]\label{thm:forcing}
Let $W\models\ZFC+V=L$ be transitive, and let $\kappa$ be an
infinite regular cardinal of $W$. Force with
\[
 \operatorname{Add}(\kappa,1)^W
 =\{p:p\text{ is a partial map }\kappa\to2,
                    \ |\dom(p)|^W<\kappa\},
\]
ordered by extension, and let $V=W[G]$. Then
\begin{enumerate}[label=\textup{(\roman*)}]
\item $\HOD^V=W$ and $\KH^V=\No^W$ as ordered fields;
\item the least birthday of a surreal outside $\KH^V$ is exactly
$\lambda_H^V=\kappa$;
\item $\Cell^V$ contains the explicit ordinal-indexed independent
family of Theorem~\ref{thm:dichotomy} over the old field $\No^W$;
\item if $\kappa>\omega$, then $\R^V=\R^W$, and every real of $V$
is ordinal definable in $V$, despite \textup{(ii)} and \textup{(iii)}.
\end{enumerate}
\end{theorem}
\begin{proof}
For two conditions $p,q$, flip the bits on the set of common-domain
coordinates at which they disagree. This is a ground-model
automorphism sending $p$ to a condition compatible with $q$.
Thus the forcing is weakly homogeneous and ordinal definable, so
Lemma~\ref{lem:forcingHOD} applies. The field identification follows
from Theorem~\ref{thm:HOD} and the absoluteness of canonical surreal
arithmetic.

Regularity of $\kappa$ makes the union of a descending chain of
length less than $\kappa$ a condition. Given a name for a sequence
of ordinals of length $\mu<\kappa$, recursively strengthen a
condition to decide each coordinate, taking lower bounds at limit
stages and at the end. The resulting ground-model sequence is
forced to equal the name. Hence the forcing adds no subsets of any
$\mu<\kappa$. For $\kappa=\omega$, the corresponding assertion
about finite subsets is immediate. This closure also prevents a
collapse of the cardinal $\kappa$ by a shorter sequence.

The union $g=\bigcup G$ is a total subset of $\kappa$, and is new:
for every ground-model $a\subseteq\kappa$, the set of conditions
already disagreeing with $a$ at some coordinate is dense. Since
$\HOD^V=W$, the first non-HOD subset therefore occurs exactly at
$\kappa$. Theorem~\ref{thm:firstbirthday} proves (ii).
The new sign sequence with positive positions $g$ lies outside
$\KH^V$; normalize its absolute value as in
Theorem~\ref{thm:amplify} to obtain the parameter $t$ for (iii).

When $\kappa>\omega$, the no-new-short-sequences conclusion in
particular gives no new subsets of $\omega$, hence no new reals.
All old reals belong to $W=\HOD^V$, which proves (iv).
\end{proof}

\begin{example}[All reals ordinal definable, but many omnific integers not]
Take $\kappa=\omega_1^W$ in Theorem~\ref{thm:forcing}. In the
extension, every real is ordinal definable; every surreal of birthday
less than $\omega_1^W$ is in the HOD-surreal field; and a non-HOD
surreal first appears at birthday $\omega_1^W$. Nevertheless a
single fixed interval of purely infinite omnific integers already
contains independent sets of every internal cardinality over the
old surreal field. This is a strict separation between the real
question and the surreal question, not a consequence of adding an
undefinable real.
\end{example}

\begin{remark}[What is, and is not, absolute]
The forcing theorem compares old sign sequences and their arithmetic.
It does not assert that every old parameter-free definition keeps
its truth or uniqueness in an extension. An ambient definition can
mention the continuum, a generic object, or the set-theoretic
universe itself. Its truth can change even when the particular sign
sequence and all arithmetic performed on old inputs stay unchanged.
Nor is the relation $V=\HOD$ preserved by arbitrary forcing.
\end{remark}

\section{Relativization and a surcomplex companion}
\label{sec:relativize}

\subsection{An allowed set of ordinal information}

Let $p$ be a fixed set of ordinals. Write $\OD(p)$ for definitions
using $p$ together with finitely many ordinals, and $\HOD(p)$ for
hereditary membership in this class. Here the parameter is the
\emph{single set} $p$, not permission to use its arbitrary elements
as unrelated parameters. Since $p$ is a set of ordinals, it belongs
to $\HOD(p)$; this is an inner model of ZFC with its corresponding
canonical definable well-order \cite{Jech}.

\begin{proposition}[Relativized theory]\label{prop:relative}
Put $K_p=\No\cap\OD(p)$ and $I_p=\Oz\cap\OD(p)$. Then
\[
 K_p=\No\cap\HOD(p)=\No^{\HOD(p)},\qquad
 K_p=\Frac(I_p).
\]
The field $K_p$ is real closed and $I_p$ is an integer part.
Every member of $\Cell$ is in $\OD(p)$ if and only if
$V=\HOD(p)$. If $V\ne\HOD(p)$, the family
\[
 \omega+\omega^{\,t\omega^{-(\alpha+1)}}\quad(\alpha\in\Ord)
\]
is algebraically independent over $K_p$ for every
$t\in(0,1)\setminus K_p$.
\end{proposition}
\begin{proof}
The hereditary sign-code argument only uses that ordinals and the
fixed sign tags are hereditarily definable; permitting $p$ does not
change it. Unique-output closure gives the field, floor, and fraction
claims. The reversible code uses no extra parameter, while the
membership-graph collapse works inside $\HOD(p)$ because that inner
model contains $p$. Finally every normal-form entry of an
$\OD(p)$ surreal is definable from $p$ and ordinal parameters, so
its support lies in $K_p$. The support-subfield amplification proof
applies without change.
\end{proof}

This relativization describes precisely how adding allowed information
changes the theory. It is different from allowing \emph{all} surreal
parameters, which makes every individual surreal trivially definable
by equality to itself.

\subsection{Canonical surcomplex pairs}

For completeness, define the surcomplex field by pairs
$\No(i)=\No\times\No$ with $i=(0,1)$ and the usual quadratic
extension operations. This fixed coordinate representation names a
real axis and an imaginary unit.

\begin{proposition}[Definable surcomplex closure]
\label{prop:surcomplex}
In the canonical pair representation, the ambiently $A$-definable
surcomplex numbers are exactly $K_A(i)$, and they form an algebraically
closed field. The ordinal-definable ones are exactly $\KH(i)$.
If $V\ne\HOD$, the real omnific family in
Theorem~\ref{thm:dichotomy} remains algebraically independent over
$\KH(i)$ inside $\No(i)$.
\end{proposition}
\begin{proof}
Pairing and the two coordinate projections are parameter-free
set-theoretic operations. Thus a pair is $A$-definable precisely when
both coordinates are. The quadratic extension of a real closed field
by $i$ is algebraically closed. For the independence assertion,
write a purported complex-coefficient polynomial relation as
$P_0(\bar z)+iP_1(\bar z)=0$, with polynomials over $\KH$.
At real inputs both values are real, so both are zero; a nonzero
relation would contradict real algebraic independence.
\end{proof}

There is no conflict with complex conjugation or with the failure of
a pure field to choose an imaginary unit. The proposition uses
\emph{ambient canonical coordinates}. An abstract algebraically
closed-field language without those coordinates has a different
definable closure, just as the abstract real field does.

\section{Further questions and a formalization route}
\label{sec:questions}

The preceding theorems have complete written proofs. The questions
below are further projects, not omitted hypotheses or unfinished
steps in those proofs. They are not asserted to be previously
published open problems.

\begin{question}[The omnific predicate as a language expansion]
Determine the parameter-free definable closure of
$(\No;+,-,\cdot,<,0,1,\Oz)$, or of well-specified set-sized models
of its first-order theory. Which ambiently definable surreals become
intrinsically definable after adding only the omnific predicate?
\end{question}
The pure ordered-field answer is Theorem~\ref{thm:purefield}.
The omnific code does not settle the expanded-language question:
its decoder recognizes ordinal-indexed monomials and normal-form
coefficients through set theory, not just through the ring predicate.

\begin{question}[Birthday cost of reversible omnific coding]
Obtain sharp upper and lower bounds for $\bd(\Code(s))$ in terms of
$\bd(s)$, and compare alternative reversible codes with the same
leading term and coefficient alphabet. Can a useful optimality
statement be proved for a specified class of such codes?
\end{question}
Our construction controls support order type and the leading scale,
not the exact surreal birthday of its nested monomials. No unstated
birthday estimate is used in the independence theorem.

\begin{question}[Intermediate support fields]
Classify naturally occurring proper subfields satisfying
\eqref{eq:Khyp}, especially those arising from relative definability
or canonical inner models. Determine when stronger structural
invariants than the independent family can be read from the
quotient additive group $\No/K$.
\end{question}
The support-subfield theorem already gives one invariant uniformly:
a single element outside such a field produces an ordinal-indexed
independent family in one fixed omnific interval.

\subsection{A realistic Lean decomposition}

A formalization should separate the general Hahn-field mathematics
from satisfaction and internal models. The shortest algebraic route
is to formalize Lemma~\ref{lem:cosets} for set-sized Hahn series and
finite polynomials first: coefficient supports in different cosets
cannot cancel. Next formalize the finite linear independence of
monomials, the normalized-parameter argument, and the finite form
of Theorem~\ref{thm:amplify}. These steps require no model-theoretic
truth predicates or proper-class objects.

The set-theoretic track then needs a precise universe convention,
canonical sign codes, unique normal forms, and definable class maps.
HOD identification requires a model-of-ZFC interface, with standard
external formulas kept distinct from internal formula codes.
The statements about $B_n$ should be separate theorems for each
fixed fragment, not a falsely uniform truth evaluator.
The repository already documents sign, Hahn, size, and normal-form
interfaces, but this article neither maps the new statements to
specific checked declarations nor claims to have built or extended
its Lean code \cite{Repo}.

\section{Conclusion}

There is no single language-independent class of ``the definable
surreals.'' Once the ambient universe and permitted parameters are
fixed, however, the basic algebra is robust: definable surreals form
a real closed field, and definable omnific integers are an integer
part with exactly that fraction field. Ordinal parameters reveal a
particularly natural inner field, the canonical surreal field of HOD.

The strongest results are not counting observations. An explicit
normal-form map compresses every surreal definability degree into a
fixed-leading-term omnific interval. If HOD fails to contain the
universe, one missing surreal generates an ordinal-indexed
algebraically independent family in that same interval. Uniform
coefficient assignments, rather than small coefficient alphabets or
ordinary congruences, carry the missing information. This yields a
concrete bridge from set-theoretic definability to the support
geometry and transcendence theory of surreal and omnific arithmetic.

\appendix
\section{Provenance, proof scope, and reproducibility}
\label{app:audit}

\subsection{Claim ledger}

\begin{longtable}{@{}P{0.26\textwidth}P{0.29\textwidth}P{0.39\textwidth}@{}}
\toprule
Statement family & Status & Dependency and qualification\\
\midrule
\endhead
Signs, normal forms, real closedness, and omnific floor & Classical &
Conway, Gonshor, van den Dries--Ehrlich, and
L'Innocente--Mantova; elementary proofs supplied where useful.\\[4pt]
Ambient field and integer-part closure & Standard-method consequence &
Unique-output closure, ordered finite root selection, and a uniformly
chosen ordinal denominator. No priority claim for the general
closure principle.\\[4pt]
Global birthday/set-theory interpretation & Prior announced work &
Chen--Hamkins--Yang. Not used as a black-box proved dependency and
not claimed here.\\[4pt]
HOD/sign-code identification & Foundational application &
Hereditariness of a sign code, inner-model recursion, and canonical
normal forms. No claim to have introduced HOD-valued surreal fields.\\[4pt]
Fixed-interval reversible code & Proposed contribution &
Theorem~\ref{thm:code}; complete normal-form proof. Search did not
locate this precise fixed-leading-term code.\\[4pt]
One-interval universe tests & Proposed consequences &
Theorems~\ref{thm:globalOD} and \ref{thm:pointwise}; standard pointed
set coding followed by the new omnific code.\\[4pt]
Support-subfield amplification and HOD transcendence dichotomy &
Principal proposed contribution &
Theorems~\ref{thm:amplify} and \ref{thm:dichotomy}; finite support-coset
proof, explicit ordinal-indexed family. No transcendence-basis claim.\\[4pt]
First non-HOD birthday & Translation/refinement &
The threshold is the standard first subset disagreement
$\Delta(\HOD,V)$. Its birthday interpretation is explicit; the
compactness restriction is credited to \cite{BCHGP2026}.\\[4pt]
Bounded-complexity escape & Standard-method refinement &
Fixed-fragment truth and Replacement; no all-level truth predicate
and no asserted one-step complexity bound.\\[4pt]
Forcing without new reals & Application &
Standard homogeneous closed forcing over $L$, combined with the
new independent-family theorem.\\[4pt]
Finite verification & Regression tests only &
Finite symbolic codes, decoder validation, and formal support-coset
checks. These do not prove any infinite or set-theoretic theorem.\\
\bottomrule
\end{longtable}

\subsection{What was actually inspected}

The repository was accessed through its GitHub connector at the
pinned commit printed in Section~\ref{sec:intro}. Inspection covered
its README, the report catalogue, the introduction and provenance
section of the birthday-cutoff/hereditary-set report, and the
introduction and convention section of the omnific-Diophantine
report. It was not a line-by-line audit of all 51 reports or their
preserved source manuscripts. No existing repo files were edited,
no commit was made, and no Lean build was performed.

The literature comparison used the supplied definable-real reference,
primary works and author announcements on pointwise definability,
HOD, surreal arithmetic, omnific normal forms, and quantifier
elimination. The recent first-HOD-disagreement preprint was checked
as well. The cited works and the repository catalogue did not reveal
the specific fixed-interval code or the support-subfield
transcendence amplification proved here. This limited negative
search evidence cannot establish publication priority.

\subsection{Files and checks}

The archive contains this article as both \repopath{.tex} and
\repopath{.pdf}, a build/readme file, a standalone Python regression
script, its JSON results, and a concise audit note. The source is
self-contained apart from ordinary LaTeX packages: its bibliography
is included below, and no external graphics or bibliography database
is required. Build with \repopath{latexmk -pdf} or by running
\repopath{pdflatex} sufficiently many times to resolve references.

The finite script represents the nested exponents $e_n$ by formal
tags. It checks the decoder on all sign words through a stated
finite length, rejects malformed codes, and verifies support-coset
separation on finite formal examples. It deliberately does not
replace $\omega$ by a large real number: that numerical substitution
would destroy the exact non-Archimedean ordering relevant to the
proofs. The script neither decides definability nor approximates HOD.

The delivered run passed 8,191 finite code round trips through sign
length 12, 11 malformed-code rejections, 3,121 exact rational
normalization checks, 2,157 formal infinitesimal floor cases, and
500 finite support-coset polynomial checks. A deliberately dependent
exponent example correctly produced a vanishing polynomial, checking
that the independence hypothesis is not silently omitted. The JSON
file records the detailed counts and the fixed random seed.


\clearpage
\section{Notation and reading guide}

\begin{longtable}{@{}P{0.22\textwidth}P{0.72\textwidth}@{}}
\toprule
Notation & Meaning\\
\midrule
\endhead
$\No$, $\Oz$ & Canonical surreal field and classical omnific-integer ring.\\[3pt]
$\bd(s)$ & Ordinal length of the canonical sign sequence.\\[3pt]
$\omega^x$ & Conway monomial map; not general surreal exponential
power with base $\omega$.\\[3pt]
$\NF(x)$, $\supp(x)$ & Whole indexed normal form; its exponent support.\\[3pt]
$D_A^M$ & External definable closure of $A$ in the membership structure
$(M,\in)$, using standard formulas and finite parameter tuples.\\[3pt]
$K_A$, $I_A$ & Ambiently $A$-definable surreals and omnific integers.\\[3pt]
$\KH$, $\IH$ & Ordinal-definable surreals and omnific integers;
$\KH=\No^{\HOD}$ under canonical sign representation.\\[3pt]
$\Cell$ & $\{z\in\Oz:\omega<z<\omega+\omega^{1/2},\ \ct(z)=0\}$.\\[3pt]
$e_\alpha$ & $\omega^{-(\alpha+1)}$, a positive infinitesimal surreal.\\[3pt]
$\Code$, $\Decode$ & Reversible omnific coding map and its partial inverse.\\[3pt]
$\lambda_H$ & First ordinal with a non-HOD subset; equivalently the
first birthday of a non-OD surreal.\\[3pt]
$B_n$, $d_n$ & Fixed-$\Sigma_n$ uniquely definable surreals, and the
corresponding definable ordinal surreal escaping that level.\\
\bottomrule
\end{longtable}

For the main new algebraic argument alone, read
Sections~\ref{sec:foundations}, \ref{sec:HOD}, and
\ref{sec:transcendence}. For the definability concept and its
model-theoretic limitations, read Sections~\ref{sec:definitions},
\ref{sec:global}, and \ref{sec:complexity}. The code in
Section~\ref{sec:code} connects the two routes.

\begin{thebibliography}{99}
\addcontentsline{toc}{section}{References}

\bibitem{Conway}
J. H. Conway, \emph{On Numbers and Games}, second edition,
A K Peters, 2001; first edition, Academic Press, 1976.

\bibitem{Gonshor}
H. Gonshor, \emph{An Introduction to the Theory of Surreal Numbers},
London Mathematical Society Lecture Note Series 110,
Cambridge University Press, 1986.

\bibitem{vdDE}
L. van den Dries and P. Ehrlich,
Fields of surreal numbers and exponentiation,
\emph{Fundamenta Mathematicae} \textbf{167} (2001), 173--188;
erratum, \textbf{168} (2001), 295--297.
\href{https://doi.org/10.4064/fm167-2-3}{doi:10.4064/fm167-2-3};
\href{https://doi.org/10.4064/fm168-3-5}{erratum doi:10.4064/fm168-3-5}.
Only classical background, not a disputed cutoff formula, is used here.

\bibitem{LM}
S. L'Innocente and V. Mantova,
\emph{A factorisation theory for generalised power series and omnific
integers}, arXiv:1710.07304, version 5, January 22, 2024.
\url{https://arxiv.org/abs/1710.07304}.
Journal version: \href{https://doi.org/10.1016/j.aim.2024.109513}
{doi:10.1016/j.aim.2024.109513}.

\bibitem{Marker}
D. Marker, \emph{Model Theory: An Introduction}, Graduate Texts in
Mathematics 217, Springer, 2002.
For the quantifier-elimination argument, see also the author's
\emph{Introductory Lectures on Model Theory II: Quantifier Elimination},
October 14, 2020,
\url{https://homepages.math.uic.edu/~marker/MSRI-2.pdf}.

\bibitem{Jech}
T. Jech, \emph{Set Theory: The Third Millennium Edition, Revised and
Expanded}, Springer Monographs in Mathematics, Springer, 2003.

\bibitem{Shepherdson}
J. C. Shepherdson,
A non-standard model for a free variable fragment of number theory,
\emph{Bulletin de l'Acad\'emie Polonaise des Sciences, S\'erie des
sciences math\'ematiques, astronomiques et physiques}
\textbf{12} (1964), 79--86.
The integer-part application to $\Oz$ is also stated explicitly
in \cite[Section 1.1]{LM}.

\bibitem{HLR}
J. D. Hamkins, D. Linetsky, and J. Reitz,
Pointwise definable models of set theory,
\emph{Journal of Symbolic Logic} \textbf{78} (2013), no.~1, 139--156.
\url{https://arxiv.org/abs/1105.4597}.
Author's discussion:
\url{https://jdh.hamkins.org/pointwisedefinablemodelsofsettheory/}.

\bibitem{HamkinsLeahy}
J. D. Hamkins and C. Leahy,
Algebraicity and implicit definability in set theory,
\emph{Notre Dame Journal of Formal Logic} \textbf{57} (2016),
no.~3, 431--439.
\url{https://arxiv.org/abs/1305.5953};
\href{https://doi.org/10.1215/00294527-3542326}
{doi:10.1215/00294527-3542326}.

\bibitem{GoldbergPoveda}
G. Goldberg and A. Poveda,
Compactness phenomena in HOD,
\emph{Forum of Mathematics, Sigma} \textbf{13} (2025), e118.
\href{https://doi.org/10.1017/fms.2025.10070}
{doi:10.1017/fms.2025.10070}.
Author's version:
\url{https://math.berkeley.edu/~goldberg/Papers/CompactnessPhenomenaInHOD.pdf}.

\bibitem{BCHGP2026}
T. Benhamou, J. Cummings, G. Goldberg, Y. Hayut, and A. Poveda,
\emph{Compactness phenomena in HOD and the Optimality of Magidor's
Covering theorem}, arXiv:2608.24190, version 1, August 25, 2026.
\url{https://arxiv.org/abs/2608.24190}.
Theorems 1.1--1.2 supply the cited first-disagreement results;
these are external inputs, not newly proved claims of this article.

\bibitem{CHY}
J. D. Hamkins, \emph{Surreal arithmetic is bi-interpretable with set
theory}, CUNY Logic Workshop announcement, March 2026;
joint work with Junhong Chen and Ruizhi Yang, described as work in progress.
\url{https://jdh.hamkins.org/surreal-arithmetic-cuny-logic-workshop-march-2026/}.
The announcement, rather than a refereed completed paper, is the
source used for the attribution.

\bibitem{WikiDef}
Wikipedia contributors, \emph{Definable real number},
user-specified starting reference, accessed September 23, 2026.
\url{https://en.wikipedia.org/wiki/Definable_real_number}.
No new theorem depends on this secondary reference.

\bibitem{Repo}
V. Reshetnikov and contributors, \emph{Surreal}, GitHub repository,
README and report catalogue at commit \repoSHA,
accessed September 23, 2026.
\url{https://github.com/VladimirReshetnikov/Surreal/tree/9a385d3957bdfe3d9ea79f9a524751c90bd2c894}.

\bibitem{RepoBirthday}
\emph{Birthday Cutoffs and Hereditary Sets}, AI-assisted merged research
report in the Surreal repository, pinned snapshot as in \cite{Repo}.
\repopath{docs/foundations-and-computation/birthday-cutoffs-and-hereditary-sets/article.tex}.
The introduction explicitly credits Chen--Hamkins--Yang's announced
global bi-interpretation and separates that work from cutoff refinements.

\bibitem{RepoOmnific}
\emph{Omnific Integers and Omnific--Diophantine Geometry}, AI-assisted
merged research report in the Surreal repository, pinned snapshot
as in \cite{Repo}.
\repopath{docs/surreal/omnific-diophantine-geometry/article.tex}.
The definition and convention sections were inspected; nonclassical
Diophantine claims are not dependencies of the present proofs.

\end{thebibliography}
\end{document}
